\documentclass[11pt]{article}

% ---- Packages (all standard; compiles with a normal pdflatex install) ----
\usepackage[utf8]{inputenc}
\usepackage[T1]{fontenc}
\usepackage{lmodern}
\usepackage[a4paper,margin=1in]{geometry}
\usepackage{xcolor}
\usepackage{listings}
\usepackage{booktabs}
\usepackage{enumitem}
\usepackage{titlesec}
\usepackage{fancyhdr}
\usepackage{parskip}
\usepackage[hidelinks]{hyperref}

% ---- Colours ----
\definecolor{accent}{HTML}{1d4ed8}
\definecolor{codebg}{HTML}{0d1117}
\definecolor{codetext}{HTML}{e6edf3}
\definecolor{codekw}{HTML}{ff7b72}
\definecolor{codestr}{HTML}{a5d6ff}
\definecolor{codecom}{HTML}{8b949e}

\hypersetup{colorlinks=true, linkcolor=accent, urlcolor=accent, citecolor=accent}

% ---- Headers / footers ----
\pagestyle{fancy}
\fancyhf{}
\renewcommand{\headrulewidth}{0.4pt}
\lhead{\small\itshape VoteChain: Decentralised Voting System}
\rhead{\small\thepage}
\cfoot{\small Department of Physics, IIT Kanpur}

% ---- Section styling ----
\titleformat{\section}{\Large\bfseries\color{accent}}{\thesection.}{0.6em}{}
\titleformat{\subsection}{\large\bfseries}{\thesubsection}{0.6em}{}

% ---- Code listing style ----
\lstdefinestyle{code}{
  backgroundcolor=\color{codebg},
  basicstyle=\ttfamily\footnotesize\color{codetext},
  keywordstyle=\color{codekw},
  stringstyle=\color{codestr},
  commentstyle=\color{codecom}\itshape,
  numberstyle=\tiny\color{codecom},
  numbers=left, numbersep=8pt,
  breaklines=true, frame=single, framerule=0pt,
  rulecolor=\color{codebg},
  showstringspaces=false, tabsize=2,
  xleftmargin=10pt, xrightmargin=6pt, aboveskip=10pt, belowskip=6pt,
}
\lstset{style=code}
\lstdefinelanguage{Solidity}{
  keywords={pragma,solidity,contract,function,public,private,view,returns,mapping,struct,enum,address,uint,uint256,bool,string,bytes32,memory,storage,require,emit,event,modifier,constructor,if,else,for,return,msg,keccak256,abi},
  keywordstyle=\color{codekw},
  comment=[l]{//}, morecomment=[s]{/*}{*/},
  morestring=[b]",
  sensitive=true,
}

\begin{document}

% ================= TITLE =================
\begin{titlepage}
\centering
\vspace*{2cm}
{\Huge\bfseries VoteChain\par}
\vspace{0.5cm}
{\LARGE A Decentralised Voting System Using\\[0.2cm] Ethereum Smart Contracts\par}
\vspace{1.5cm}
{\large\bfseries Yash Yadav\par}
\vspace{0.3cm}
{\large Department of Physics, Indian Institute of Technology Kanpur\par}
\vspace{0.2cm}
{\large 2024--25\par}
\vspace{2cm}
\begin{minipage}{0.85\textwidth}
\small
\textbf{Abstract.} Traditional electronic voting systems suffer from single points of
failure, opaque audit trails, and centralised control that requires trust in an
authority. This project presents \emph{VoteChain}, a full-stack decentralised voting
platform built on the Ethereum blockchain. The system enforces election integrity
through immutable smart-contract logic written in Solidity, role-based access control
separating super-admin, election-admin, and voter roles, and an event-sourced frontend
that reads on-chain data gas-free. Two voting modes are supported: a transparent
one-vote-per-wallet ballot and a privacy-preserving \emph{commit--reveal} scheme
implemented fully on-chain, in which votes are submitted as cryptographic commitments
and tallied only after a separate reveal phase. Administrative setup is batched so that
an entire candidate slate or voter roll is registered in a single transaction. The
contract is deployed on the Ethereum Sepolia testnet and is validated by a 48-case
Hardhat test suite covering access control, the election lifecycle, double-vote
prevention, and the commit--reveal flow. The report also documents the security
engineering applied to the repository itself, including secret-management remediation.
\end{minipage}
\vfill
{\small Live contract: \url{https://sepolia.etherscan.io/address/0xBD29A19DBD25bD4d97C09b65862CB7061A9b3011}\par}
\end{titlepage}

\tableofcontents
\newpage

% ================= 1. INTRODUCTION =================
\section{Introduction}

Elections are foundational to democratic systems, yet existing electronic voting
infrastructure frequently concentrates control in a single authority, making
comprehensive external audit difficult. Paper-based systems are slow and prone to
physical manipulation. Blockchain-based voting addresses the integrity and auditability
dimensions of this problem by recording votes as cryptographically signed on-chain
transactions that are visible to all participants and modifiable by none.

\subsection{Motivation}
Recent public discourse around election integrity motivates a system in which:
\begin{itemize}[nosep]
  \item Every vote is an immutable on-chain transaction.
  \item Counting is deterministic and publicly verifiable.
  \item No central authority controls data after submission.
  \item The complete audit trail is available to any observer.
\end{itemize}

\subsection{Scope and Limitations}
Blockchain guarantees \emph{integrity} (votes cannot be modified) and \emph{transparency}
(totals are publicly verifiable). It does not, by itself, guarantee voter \emph{identity}
(a wallet is not a verified person). Vote \emph{privacy} on a public chain is addressed in
this project by an on-chain commit--reveal protocol. Identity verification is delegated to
an off-chain admin whitelist: eligible voters register their public wallet address with the
election authority, which authorises those addresses on-chain. The authentication primitive
is therefore wallet ownership --- a voter proves control of a registered address by signing
with its private key; nothing is ever typed as a password.

% ================= 2. BACKGROUND =================
\section{Background}

\subsection{Ethereum and Smart Contracts}
Ethereum is a decentralised, Turing-complete blockchain on which programs called
\emph{smart contracts} execute deterministically across all network nodes. A contract
deployed to an address is immutable; its logic cannot be changed after deployment.
Transactions that call contract functions are cryptographically signed by the sender's
private key, creating a tamper-evident audit trail. Gas fees denominated in Ether
compensate validators for computation.

\subsection{Solidity}
Solidity is the primary high-level language for Ethereum smart contracts. Features used in
this project include \emph{mappings} (hash-map storage), \emph{enums} (the election state
machine), \emph{events} (cheap on-chain logging queried by the frontend), \emph{modifiers}
(access-control predicates), and \texttt{keccak256} hashing (the commit--reveal scheme).

\subsection{MetaMask and Ethers.js}
MetaMask is a browser extension that manages Ethereum key pairs and signs transactions on
behalf of the user. Ethers.js v6 is a JavaScript library that provides a high-level API for
connecting to nodes, instantiating contracts, and querying event logs without gas cost.

% ================= 3. ARCHITECTURE =================
\section{System Architecture}

\subsection{Overview}
All persistent state lives on-chain. The frontend is a pure read/write client with no
backend server. Read operations (fetching election lists, candidates, voter rolls, vote
counts) query the contract's event logs --- a gas-free operation using the Ethereum
\texttt{eth\_getLogs} JSON-RPC call. Write operations (creating elections, casting votes,
etc.) send signed transactions via MetaMask.

\begin{quote}
Voter/Admin Browser (React + Ethers.js + MetaMask) $\rightarrow$ Ethereum RPC (Alchemy /
public Sepolia) $\rightarrow$ Deployed \texttt{VotingSystem} contract (Sepolia)
\end{quote}

\subsection{Component Summary}
\begin{center}
\begin{tabular}{@{}lll@{}}
\toprule
\textbf{Component} & \textbf{Technology} & \textbf{Role} \\
\midrule
Smart contract & Solidity 0.8.20 & On-chain state, access control, event log \\
Deploy script & Hardhat + Ethers & Sepolia deployment \\
Test suite & Hardhat + Chai (48 tests) & Correctness and security verification \\
Frontend & React + Vite & Admin and voter UI, wallet integration \\
Wallet & MetaMask & Key management, transaction signing \\
RPC provider & Alchemy / Sepolia & Ethereum node access \\
\bottomrule
\end{tabular}
\end{center}

% ================= 4. CONTRACT DESIGN =================
\section{Smart Contract Design}

\subsection{State and Data Structures}
Elections are stored in a \texttt{private} mapping to prevent direct external enumeration;
the frontend reconstructs the list from emitted events. Each \texttt{Election} contains
nested mappings for candidates, voters, and admin addresses, providing $\mathcal{O}(1)$
lookup. A per-election boolean \texttt{commitReveal} selects the voting mode.

\begin{lstlisting}[language=Solidity]
enum ElectionState { Created, Active, Reveal, Ended }

struct Voter {
    bool isAuthorized;
    bool hasVoted;        // plain: voted | commit-reveal: revealed
    uint votedCandidateId;
    bytes32 commitment;   // keccak256(candidateId, secret, voter)
    bool hasCommitted;
}
\end{lstlisting}

\subsection{Election Lifecycle}
Plain elections progress \texttt{Created} $\rightarrow$ \texttt{Active} $\rightarrow$
\texttt{Ended}. Commit--reveal elections insert a reveal window:
\texttt{Created} $\rightarrow$ \texttt{Active} (commit) $\rightarrow$ \texttt{Reveal}
$\rightarrow$ \texttt{Ended}. Candidates may only be added in \texttt{Created}; results may
only be read in \texttt{Ended}. The state machine prevents premature result leakage and
invalid transitions.

\subsection{Access Control}
A two-tier role model is enforced by \texttt{require} guards:
\begin{itemize}[nosep]
  \item \textbf{Super Admin} (contract deployer): creates elections, assigns election-level admins.
  \item \textbf{Election Admin}: adds candidates, authorises voters, starts/ends the election, opens the reveal phase.
\end{itemize}
These modifiers are the authoritative enforcement layer: even if the user interface were
bypassed, the contract itself rejects administrative transactions from non-admin wallets.

\subsection{Vote Integrity (Plain Mode)}
Four on-chain checks prevent vote manipulation: the election must be \texttt{Active}; the
voter must be authorised; the voter must not have already voted (the \texttt{hasVoted} flag
is set atomically within the same transaction, making double-spending impossible); and the
candidate id must be in range.

\subsection{Commit--Reveal Privacy (On-Chain)}
The privacy scheme is implemented as two phases.
\begin{enumerate}[nosep]
  \item \textbf{Commit:} during the \texttt{Active} window the voter submits
  $H = \mathrm{keccak256}(\texttt{candidateId} \,\|\, \texttt{secret} \,\|\,
  \texttt{voterAddress})$. The hash reveals nothing about the choice, and binding the
  voter's address prevents a commitment from being copied and replayed by another wallet.
  \item \textbf{Reveal:} after the admin opens the \texttt{Reveal} window, the voter submits
  the plaintext \texttt{(candidateId, secret)}. The contract recomputes $H$ and counts the
  vote only if it matches the stored commitment. Unrevealed commitments are never tallied.
\end{enumerate}

\begin{lstlisting}[language=Solidity]
function revealVote(uint _electionId, uint _candidateId, string memory _secret) public {
    // ... state and authorization checks ...
    bytes32 check = keccak256(abi.encodePacked(_candidateId, _secret, msg.sender));
    require(check == voter.commitment, "Reveal does not match commitment");
    voter.hasVoted = true;
    e.candidates[_candidateId].voteCount++;
    emit VoteRevealed(_electionId, msg.sender);
}
\end{lstlisting}

\subsection{Batched Administration}
To minimise the number of wallet signatures during setup, the contract exposes
\texttt{addCandidates(id, string[])} and \texttt{authorizeVoters(id, address[])}, which
register an entire list in a single transaction. A full election setup therefore requires
three signatures (create, add all candidates, authorise all voters) rather than one per
candidate and per voter. This is a deliberate design choice: every state change must be
signed --- there is no mechanism to skip confirmation for individual writes --- so the
correct optimisation is to reduce the \emph{number} of transactions, not to weaken the
signing requirement.

\subsection{Event Architecture}
Ten events enable gas-free data retrieval: \texttt{ElectionCreated},
\texttt{CandidateAdded}, \texttt{VoterAuthorized}, \texttt{VoteCast},
\texttt{VoteCommitted}, \texttt{VoteRevealed}, \texttt{ElectionStarted},
\texttt{RevealStarted}, \texttt{ElectionEnded}, and \texttt{AdminAssigned}. No event
parameter is declared \texttt{indexed}; the frontend fetches each event type and filters by
election id in JavaScript.

% ================= 5. FRONTEND =================
\section{Frontend Implementation}

\subsection{Architecture Decisions}
The application uses a single \texttt{view} state variable
(\texttt{"connect"} $|$ \texttt{"admin"} $|$ \texttt{"voter"}) set atomically with all auth
state, avoiding a routing race condition. All sub-components are defined at module scope so
React never remounts them on re-render, which preserves input focus across keystrokes.

\subsection{Identity and the ``Login'' Model}
A deliberate design point: \emph{the wallet is the login}. Connecting MetaMask proves the
user controls a particular address; there is no address field to type, because typing an
address would let anyone impersonate any voter. On connecting, the app reads
\texttt{superAdmin()} and routes the deployer wallet to the admin dashboard and every other
wallet to the voter interface. The voter view displays an explicit
\textbf{Authorized / Not authorized} badge for the connected wallet, and the admin portal
refuses to render for non-admin wallets (defence in depth atop the on-chain modifiers).

\subsection{Event-Querying Strategy}
Public Sepolia RPC nodes reject \texttt{eth\_getLogs} queries spanning more than
$\sim$500{,}000 blocks, so the frontend computes a safe range:
\begin{lstlisting}[language=Solidity]
const latest = await provider.getBlockNumber();
const fromBlock = Math.max(0, latest - 500000); // ~70 days
\end{lstlisting}

\subsection{Commit--Reveal in the UI}
Creating an election offers a ``Private (commit--reveal)'' toggle. In a private election the
voter first commits a hashed vote (the secret is auto-generated if left blank and saved in
\texttt{localStorage} for convenience), the admin then opens the reveal phase, and the voter
reveals to have the vote counted. The reveal form can pre-fill the saved commitment or accept
a manually entered candidate and secret.

% ================= 6. SECURITY =================
\section{Security Analysis}

\subsection{Threat Model}
\begin{center}
\begin{tabular}{@{}lll@{}}
\toprule
\textbf{Attack} & \textbf{Protection} & \textbf{Mechanism} \\
\midrule
Vote tampering & Blockchain immutability & EVM consensus \\
Double voting & \texttt{hasVoted} flag & Set atomically with the vote \\
Unauthorised voting & \texttt{isAuthorized} check & Admin whitelist per election \\
Vote privacy & Commit--reveal & Hashed commitment, delayed reveal \\
Commitment replay & Address-bound hash & \texttt{keccak256(..., msg.sender)} \\
Fake result reporting & On-chain counting & Tally incremented by contract \\
Admin escalation & Modifier guards & \texttt{onlySuperAdmin}/\texttt{onlyElectionAdmin} \\
Sybil (fake wallets) & Off-chain identity & Admin whitelisting boundary \\
\bottomrule
\end{tabular}
\end{center}

\subsection{Repository Security Engineering}
Beyond the contract, the project's source repository was hardened. An early commit had
inadvertently tracked a \texttt{.env} file containing a deployer private key and API
credentials, along with the \texttt{node\_modules} directory. Remediation involved rotating
all exposed credentials, rewriting Git history to purge the secrets from every commit,
correcting the \texttt{.gitignore} patterns (which were root-anchored and failed to match
nested directories), and adding a committed \texttt{.env.example} template. This mirrors a
real incident-response workflow and underlines that key management is part of a voting
system's security boundary.

\subsection{Known Limitations}
A wallet address is not a verified human identity; the system assumes off-chain identity
verification before authorisation. The contract has not undergone a formal third-party
audit. A compromised admin key grants election control --- a multi-signature admin model is
identified as future work. The commit--reveal scheme is vulnerable to weak-secret dictionary
attacks if voters choose low-entropy secrets.

\subsection{Testing and Verification}
Correctness is verified by a Hardhat test suite of \textbf{48 passing tests} executed on a
local EVM. Coverage includes: deployment and role assignment; election creation and the full
state machine; single and batched candidate/voter registration; plain voting with
double-vote, authorisation, and candidate-range checks; and the complete commit--reveal flow
(commit acceptance and rejection, reveal-phase gating, commitment matching, double-reveal
prevention, and mode isolation between plain and private elections). The suite runs with
\texttt{npm test}.

\subsection{Gas Analysis}
\begin{center}
\begin{tabular}{@{}lrr@{}}
\toprule
\textbf{Operation} & \textbf{Gas (approx.)} & \textbf{Cost @ 10 Gwei} \\
\midrule
Deploy contract & $\sim$1{,}300{,}000 & $\sim$0.013 ETH \\
createElection & $\sim$120{,}000 & $\sim$0.0012 ETH \\
addCandidates (per name) & $\sim$75{,}000 & $\sim$0.0008 ETH \\
authorizeVoters (per voter) & $\sim$45{,}000 & $\sim$0.0005 ETH \\
vote / commitVote / revealVote & $\sim$55{,}000 & $\sim$0.0006 ETH \\
Read operations / event queries & 0 & Free \\
\bottomrule
\end{tabular}
\end{center}
All read operations are view-function calls or event-log queries and are gas-free. Only
state-mutating transactions incur gas; voters pay only their own \texttt{vote}/\texttt{commit}/
\texttt{reveal} gas, while administrative gas is borne by the admin wallet.

% ================= 7. DEPLOYMENT =================
\section{Deployment}

\subsection{Technology Stack}
\begin{center}
\begin{tabular}{@{}ll@{}}
\toprule
\textbf{Layer} & \textbf{Technology / Version} \\
\midrule
Smart contract language & Solidity 0.8.20 \\
Contract framework & Hardhat 2.22.x \\
Test framework & Hardhat + Chai matchers (ethers v5) \\
Blockchain & Ethereum Sepolia testnet \\
Frontend framework & React 19 + Vite \\
Blockchain library & Ethers.js v6 (frontend) \\
Wallet & MetaMask \\
\bottomrule
\end{tabular}
\end{center}

\subsection{Deployment Process}
\begin{enumerate}[nosep]
  \item \textbf{Environment.} Configure a root \texttt{.env} with the Sepolia RPC URL and deployer private key (the Hardhat config resolves this file from the repository root).
  \item \textbf{Compile.} \texttt{npx hardhat compile}.
  \item \textbf{Test.} \texttt{npx hardhat test} (48 passing).
  \item \textbf{Deploy.} \texttt{npx hardhat run scripts/deploy.js --network sepolia}.
  \item \textbf{Configure.} Update \texttt{CONTRACT\_ADDRESS} in \texttt{contract.js} with the deployed address.
  \item \textbf{Frontend.} \texttt{cd voting-frontend \&\& npm run dev}.
\end{enumerate}
Live contract:
\url{https://sepolia.etherscan.io/address/0xBD29A19DBD25bD4d97C09b65862CB7061A9b3011}

% ================= 8. RESULTS =================
\section{Results and Demo Flow}
A complete end-to-end demonstration proceeds as follows:
\begin{enumerate}[nosep]
  \item Admin connects the deployer wallet and is auto-routed to the admin dashboard.
  \item Admin creates an election (optionally toggling private commit--reveal mode).
  \item Admin adds the entire candidate slate in one transaction.
  \item Admin authorises the full voter roll in one transaction.
  \item Admin starts the election; a live indicator appears.
  \item Each voter connects \emph{their own} registered wallet; the authorisation badge confirms eligibility.
  \item Voters cast (plain) or commit (private) their votes.
  \item For private elections the admin opens the reveal phase and voters reveal.
  \item Admin ends the election; results become readable and a bar chart renders with the winner highlighted.
\end{enumerate}
All intermediate steps are verifiable on Sepolia Etherscan by inspecting the contract's
transaction history.

% ================= 9. DISCUSSION =================
\section{Discussion}
The core contribution of this architecture is the \textbf{trust model}. In a traditional
system, a server administrator can modify vote counts without detection. Here, the counting
logic is encoded in contract bytecode verified against the on-chain hash, executed
identically by all EVM nodes, and every mutation is permanently logged. An auditor needs only
to read the chain --- no privileged access to a server is required. The honest engineering
trade-offs are acknowledged throughout: blockchain guarantees integrity but not identity, and
every on-chain write has a gas cost that creates friction at scale, partially mitigated here
by batching and event-sourced reads.

% ================= 10. FUTURE WORK =================
\section{Future Work}
\begin{itemize}[nosep]
  \item \textbf{ZK-proof integration.} Zero-knowledge proofs to prove voter authorisation without revealing the chosen candidate, achieving authentication and privacy simultaneously.
  \item \textbf{Indexed events.} Redeploy with \texttt{indexed} election ids to enable efficient on-chain topic filtering.
  \item \textbf{IPFS metadata.} Store candidate photos, manifestos, and descriptions off-chain with on-chain content hashes.
  \item \textbf{Hardware-wallet key management.} Replace software key custody with Ledger/Trezor integration for high-security deployments.
  \item \textbf{Multi-signature admin.} Require $k$-of-$n$ admin signatures to start or end an election, removing the single-point-of-failure admin key.
  \item \textbf{Real identity layer.} Integrate institutional SSO or government-ID-linked wallet binding for authenticated voter registration.
  \item \textbf{Reveal-window enforcement.} Enforce the stored start/end timestamps on-chain so phases close automatically.
\end{itemize}

% ================= 11. CONCLUSION =================
\section{Conclusion}
VoteChain demonstrates that a functional, secure, and auditable voting system can be built on
Ethereum with a compact Solidity contract and a React frontend. The key design decisions ---
immutable on-chain state, event-sourced reads, role-based access control, wallet-based
authentication, batched administration, and a genuine on-chain commit--reveal privacy scheme
--- collectively produce a system where vote counts cannot be tampered with after submission,
double voting is prevented cryptographically, results are publicly verifiable without any
trusted intermediary, and the admin can change only the lifecycle of an election, never the
votes. Validated by a 48-case test suite and deployed on Sepolia, the project also surfaces
honest trade-offs: blockchain secures integrity but not identity or privacy on its own, and
the commit--reveal extension and proposed ZK and multi-signature work chart a credible path
toward real-world deployment.

% ================= REFERENCES =================
\begin{thebibliography}{9}
\bibitem{nakamoto} S. Nakamoto, ``Bitcoin: A peer-to-peer electronic cash system,'' 2008.
\bibitem{buterin} V. Buterin, ``Ethereum: A next-generation smart contract and decentralised application platform,'' 2014.
\bibitem{solidity} Solidity Documentation, \emph{Solidity v0.8.20}, Ethereum Foundation, 2024. \url{https://docs.soliditylang.org}
\bibitem{ethers} R. Rimple, \emph{Ethers.js Documentation v6}, 2024. \url{https://docs.ethers.org/v6/}
\bibitem{cramer} R. Cramer, R. Gennaro, and B. Schoenmakers, ``A secure and optimally efficient multi-authority election scheme,'' \emph{EUROCRYPT '97}, Springer LNCS, 1997.
\bibitem{benaloh} J. Benaloh and D. Tuinstra, ``Receipt-free secret-ballot elections,'' \emph{Proc. 26th ACM STOC}, 1994.
\bibitem{hardhat} Nomic Foundation, \emph{Hardhat Documentation}, 2024. \url{https://hardhat.org}
\end{thebibliography}

\end{document}
